    The initial evidence presented here supports 
    a valuable role for an LLM in a screening process. An LLM is effective in 
    pulling out those rare public announcements with collusive content from a 
    massive set of company transcripts. It is able to effectively substitute 
    for human experts in performing an initial screen. However, given the 
    high rate of false positives and the importance of a competition agency avoiding 
    the allocation of resources to false leads, human oversight of the LLM's 
    output is critical. Here is a candidate protocol:
    \begin{description}
     \item[Step 1:] LLM reads and rates a corpus of transcripts.
     \item[Step 2:] For those transcripts with the highest scores, human experts read the LLM-generated excerpts and identify those with collusive content.
     \item[Step 3:] For a transcript surviving step 2, human experts read the entire transcript and determine if it is worthy of further examination.
     \item[Step 4:] For those deemed worthy of further examination, all transcripts are collected for the company along with competitors' transcripts and read by human experts.
    \end{description}
    \noindent Based on the assessment of that body of transcripts, 
    it can be determined if an official investigation should be 
    initiated. The case study in the next section provides proof 
    of concept as an LLM-based screen with a very modest use of 
    human resources is able to identify a compelling case for 
    prosecution.

    Prior to pursuing cases based on public communications, it would 
    be prudent for a competition agency to release guidelines describing 
    the content that companies should avoid in their public communications. 
    While courts have recognized that public communications can result in 
    an unlawful agreement, exclusive reliance on public communications to 
    obtain a conviction is relatively unexplored territory. Issuing 
    guidelines lays the legal groundwork for public prosecution or private 
    litigation as well as deterring collusive conduct. As described in
     \textcite{harrington2022collusion}, a firm's public announcement 
     risks facilitating collusion when it refers to the conduct of rival 
     firms or the industry large. Our categorization of the content of 
     the LLM-identified transcripts provides the basis for guidelines 
    describing the various types of content that firms should avoid 
    in their public statements.